% GENERATED FILE. Edit canonical lessons, then run npm run generate:books.
% canonical-source-hash: fnv1a64:1764db328dab660d
% canonical-lessons: GE-C475-prahlen, GE-C475-angeben, GE-C475-angeber, GE-C475-bescheiden, GE-C475-eitel

\chapter{Angeber und Bescheidene}
\label{ch:ge-angeber-und-bescheidene}
\hlchaptermodality{475}

\begin{chapteropening}
\textbf{By the end of this chapter:} \emph{I can talk about showing off and modesty.}

Complete the last lesson of chapter 475: I can talk about showing off and modesty.
\end{chapteropening}

\section[prahlen]{prahlen — to boast, to brag}

\label{lesson:GE-C475-prahlen}



\begin{quote}
Before the new one: say the German for curiosity, then the German for speechless.
\end{quote}

\subsection*{You'll want to know: prahlen}
\textbf{prahlen} — \textquotedblleft{}to boast, to brag\textquotedblright{}.

\begin{quote}
Er prahlt gern mit seinem teuren Auto.
\end{quote}

\begin{quote}
\emph{He likes to brag about his expensive car.}
\end{quote}

\emph{Prahlen mit} + dative.

\begin{cousinweb}
The origin of \emph{prahlen} is not settled; EtymDB lists the German word without a source.
\end{cousinweb}

\begin{grammarlens}[title={What you've built}]
One more word for telling, explaining and describing at B1.
\end{grammarlens}

\subsection*{Your turn}
\begin{itemize}
\raggedright
  \item \emph{Say it:} \emph{prahlen}
  \item \emph{Say it:} \emph{prahlen}, once more
  \item \emph{Say it:} a sentence of your own with \emph{prahlen}
  \item \emph{Recall:} say the German for curiosity, then the German for speechless, then say \emph{prahlen} again
\end{itemize}

\begin{tcolorbox}[breakable,colback=teal!4,colframe=teal!35!black,title={Before you move on}]
What does prahlen mean? (\textquotedblleft{}To boast, to brag\textquotedblright{}.) Say it once more.
\end{tcolorbox}
\section[angeben]{angeben — to show off; to state}

\label{lesson:GE-C475-angeben}



\begin{quote}
Before the new one: say the German for speechless, then the German for to boast.
\end{quote}

\subsection*{You'll want to know: angeben}
\textbf{angeben} — \textquotedblleft{}to show off; to state\textquotedblright{}.

\begin{quote}
Mein Cousin gibt immer mit seinen Noten an.
\end{quote}

\begin{quote}
\emph{My cousin is forever showing off about his marks.}
\end{quote}

Separable: \emph{er gibt an}, \emph{hat angegeben}. On forms it means to state: \emph{das Alter angeben} — to state your age.

\begin{cousinweb}
\emph{angeben} is built on \emph{an} \textquotedblleft{}on\textquotedblright{} and \emph{geben} \textquotedblleft{}to give\textquotedblright{} (EtymDB), which is inherited from Old High German \emph{geban}, from Proto-Germanic \emph{gebaną}.
\end{cousinweb}

\begin{grammarlens}[title={What you've built}]
One more word for telling, explaining and describing at B1.
\end{grammarlens}

\subsection*{Your turn}
\begin{itemize}
\raggedright
  \item \emph{Say it:} \emph{angeben}
  \item \emph{Say it:} \emph{angeben}, once more
  \item \emph{Say it:} a sentence of your own with \emph{angeben}
  \item \emph{Recall:} say the German for speechless, then the German for to boast, then say \emph{angeben} again
\end{itemize}

\begin{tcolorbox}[breakable,colback=teal!4,colframe=teal!35!black,title={Before you move on}]
What does angeben mean? (\textquotedblleft{}To show off; to state\textquotedblright{}.) Say it once more.
\end{tcolorbox}
\section[der Angeber]{der Angeber — show-off}

\label{lesson:GE-C475-angeber}



\begin{quote}
Before the new one: say the German for to boast, then the German for to show off.
\end{quote}

\subsection*{You'll want to know: der Angeber}
\textbf{der Angeber} — \textquotedblleft{}show-off\textquotedblright{}.

\begin{quote}
Der neue Kollege ist ein echter Angeber.
\end{quote}

\begin{quote}
\emph{The new colleague is a real show-off.}
\end{quote}

Plural: \emph{die Angeber}. A woman is \emph{die Angeberin}.

\begin{cousinweb}
\emph{Angeber} is built on \emph{angeben} \textquotedblleft{}to show off\textquotedblright{} with \emph{-er}; \emph{geben} is inherited from Old High German \emph{geban}, from Proto-Germanic \emph{gebaną} (EtymDB).
\end{cousinweb}

\begin{grammarlens}[title={What you've built}]
One more word for telling, explaining and describing at B1.
\end{grammarlens}

\subsection*{Your turn}
\begin{itemize}
\raggedright
  \item \emph{Say it:} \emph{der Angeber}
  \item \emph{Say it:} \emph{der Angeber}, once more
  \item \emph{Say it:} a sentence of your own with \emph{der Angeber}
  \item \emph{Recall:} say the German for to boast, then the German for to show off, then say \emph{der Angeber} again
\end{itemize}

\begin{tcolorbox}[breakable,colback=teal!4,colframe=teal!35!black,title={Before you move on}]
What does der Angeber mean? (\textquotedblleft{}Show-off\textquotedblright{}.) Say it once more.
\end{tcolorbox}
\section[bescheiden]{bescheiden — modest}

\label{lesson:GE-C475-bescheiden}



\begin{quote}
Before the new one: say the German for to show off, then the German for show-off.
\end{quote}

\subsection*{You'll want to know: bescheiden}
\textbf{bescheiden} — \textquotedblleft{}modest\textquotedblright{}.

\begin{quote}
Trotz ihres Erfolgs ist die Sängerin bescheiden geblieben.
\end{quote}

\begin{quote}
\emph{Despite her success, the singer has stayed modest.}
\end{quote}

Of things it means modest, small: \emph{ein bescheidenes Gehalt}.

\begin{cousinweb}
\emph{bescheiden} is usually traced to the past participle of Middle High German \emph{bescheiden} \textquotedblleft{}to decide, to allot\textquotedblright{}, built on \emph{scheiden} \textquotedblleft{}to separate\textquotedblright{}; EtymDB lists the German word without a source.
\end{cousinweb}

\begin{grammarlens}[title={What you've built}]
One more word for telling, explaining and describing at B1.
\end{grammarlens}

\subsection*{Your turn}
\begin{itemize}
\raggedright
  \item \emph{Say it:} \emph{bescheiden}
  \item \emph{Say it:} \emph{bescheiden}, once more
  \item \emph{Say it:} a sentence of your own with \emph{bescheiden}
  \item \emph{Recall:} say the German for to show off, then the German for show-off, then say \emph{bescheiden} again
\end{itemize}

\begin{tcolorbox}[breakable,colback=teal!4,colframe=teal!35!black,title={Before you move on}]
What does bescheiden mean? (\textquotedblleft{}Modest\textquotedblright{}.) Say it once more.
\end{tcolorbox}
\section[eitel]{eitel — vain}

\label{lesson:GE-C475-eitel}



\begin{quote}
Before the new one: say the German for show-off, then the German for modest.
\end{quote}

\subsection*{You'll want to know: eitel}
\textbf{eitel} — \textquotedblleft{}vain\textquotedblright{}.

\begin{quote}
Er ist so eitel, dass er ständig in den Spiegel schaut.
\end{quote}

\begin{quote}
\emph{He is so vain that he keeps looking in the mirror.}
\end{quote}

Old meaning: pure, mere — \emph{eitel Sonnenschein}, nothing but sunshine.

\begin{cousinweb}
\emph{eitel} is inherited from Middle High German \emph{ītel} \textquotedblleft{}empty, vain\textquotedblright{} and Old High German \emph{ītal}, from Proto-Germanic \emph{īdalaz} \textquotedblleft{}empty\textquotedblright{} (EtymDB); English \emph{idle} is a cousin.
\end{cousinweb}

\begin{grammarlens}[title={What you've built}]
One more word for telling, explaining and describing at B1.
\end{grammarlens}

\subsection*{Your turn}
\begin{itemize}
\raggedright
  \item \emph{Say it:} \emph{eitel}
  \item \emph{Say it:} \emph{eitel}, once more
  \item \emph{Say it:} a sentence of your own with \emph{eitel}
  \item \emph{Recall:} say the German for show-off, then the German for modest, then say \emph{eitel} again
\end{itemize}

\begin{tcolorbox}[breakable,colback=teal!4,colframe=teal!35!black,title={Before you move on}]
What does eitel mean? (\textquotedblleft{}Vain\textquotedblright{}.) Say it once more.
\end{tcolorbox}
